\documentclass[10pt,a4paper]{amsart}
\usepackage{amsmath,amssymb,amsthm,graphicx}
\usepackage{verbatim}
\newenvironment{reference}{\comment}{\endcomment}
%\newenvironment{reference}{}{}
\newenvironment{slogan}{\comment}{\endcomment}
\newenvironment{history}{\comment}{\endcomment}

% For commutative diagrams we use Xy-pic
\usepackage[all]{xy}

% We use 2cell for 2-commutative diagrams.
\xyoption{2cell}
% Theorem environments.
%
\theoremstyle{plain}
\newtheorem{theorem}[subsection]{Theorem}
\newtheorem{proposition}[subsection]{Proposition}
\newtheorem{lemma}[subsection]{Lemma}

\theoremstyle{definition}
\newtheorem{definition}[subsection]{Definition}
\newtheorem{example}[subsection]{Example}
\newtheorem{exercise}[subsection]{Exercise}
\newtheorem{situation}[subsection]{Situation}

\theoremstyle{remark}
\newtheorem{remark}[subsection]{Remark}
\newtheorem{remarks}[subsection]{Remarks}

\numberwithin{equation}{subsection}

% Macros
%
\def\lim{\mathop{\mathrm{lim}}\nolimits}
\def\colim{\mathop{\mathrm{colim}}\nolimits}
\def\Spec{\mathop{\mathrm{Spec}}}
\def\Hom{\mathop{\mathrm{Hom}}\nolimits}
\def\Ext{\mathop{\mathrm{Ext}}\nolimits}
\def\SheafHom{\mathop{\mathcal{H}\!\mathit{om}}\nolimits}
\def\SheafExt{\mathop{\mathcal{E}\!\mathit{xt}}\nolimits}
\def\Sch{\mathit{Sch}}
\def\Mor{\mathop{\mathrm{Mor}}\nolimits}
\def\Ob{\mathop{\mathrm{Ob}}\nolimits}
\def\Sh{\mathop{\mathit{Sh}}\nolimits}
\def\NL{\mathop{N\!L}\nolimits}
\def\CH{\mathop{\mathrm{CH}}\nolimits}
\def\proetale{{pro\text{-}\acute{e}tale}}
\def\etale{{\acute{e}tale}}
\def\QCoh{\mathit{QCoh}}
\def\Ker{\mathop{\mathrm{Ker}}}
\def\Im{\mathop{\mathrm{Im}}}
\def\Coker{\mathop{\mathrm{Coker}}}
\def\Coim{\mathop{\mathrm{Coim}}}

% Boxtimes
%
\DeclareMathSymbol{\boxtimes}{\mathbin}{AMSa}{"02}

%
% Macros for moduli stacks/spaces
%
\def\QCohstack{\mathcal{QC}\!\mathit{oh}}
\def\Cohstack{\mathcal{C}\!\mathit{oh}}
\def\Spacesstack{\mathcal{S}\!\mathit{paces}}
\def\Quotfunctor{\mathrm{Quot}}
\def\Hilbfunctor{\mathrm{Hilb}}
\def\Curvesstack{\mathcal{C}\!\mathit{urves}}
\def\Polarizedstack{\mathcal{P}\!\mathit{olarized}}
\def\Complexesstack{\mathcal{C}\!\mathit{omplexes}}
% \Pic is the operator that assigns to X its picard group, usage \Pic(X)
% \Picardstack_{X/B} denotes the Picard stack of X over B
% \Picardfunctor_{X/B} denotes the Picard functor of X over B
\def\Pic{\mathop{\mathrm{Pic}}\nolimits}
\def\Picardstack{\mathcal{P}\!\mathit{ic}}
\def\Picardfunctor{\mathrm{Pic}}
\def\Deformationcategory{\mathcal{D}\!\mathit{ef}}
\usepackage{cleveref}
\crefname{originalreference}{original reference}{original references}
\Crefname{originalreference}{Original reference}{Original references}
\makeatletter
\newlabel{lemma-RS-fiber-square}{{06J3}{source}{}{sourcepointer-1}{}}
\newlabel{lemma-RS-fiber-square@cref}{{[originalreference][0][]06J3}{[1][0][]source}{}{}{}}
\makeatother
\title{The 2-fibre product of cofibred groupoids satisfying (RS) satisfies (RS)}
\author{The Stacks Project Authors}
\date{stacks-project:formal-defos.tex:local-snapshot-20261012; GNU FDL 1.2 or later}
\begin{document}
\maketitle
\noindent\textit{Editorial note.} All delivered mathematical text is the author's own, copied from the pinned Stacks Project chapter file formal-defos.tex; no statement or proof step was rewritten, reordered or removed.\par\medskip
\noindent\textit{Editorial note.} The counted proof is the author's complete argument: the translation of objects and diagrams of the 2-fibre product into triples, the construction of the candidate fibre product from the outer fibre products with the comparison morphism moved along the fibre squares in the middle category, and the verification of the fibre square condition.\par\medskip
\noindent\textit{Editorial note.} The label lemma-RS-fiber-product-morphisms names the counted statement, and the proof environment immediately following it in the source is delivered verbatim as its proof.\par\medskip
\noindent\textit{Editorial note.} The counted statement is carried by the author's own lemma environment in the chapter source.\par\medskip
\noindent\textit{Editorial note.} Pre-open gate: python operations/agent-authoring/tools/precard\_gate.py --source formal-defos.tex --label lemma-RS-fiber-product-morphisms --start 5371 --end 5441 --local-id 28023 returned RESULT=OK (L28023 free, occupancy FREE, public interval clear). Candidate taken from sources/stacks-formal-defos/candidates-halfA.json.\par\medskip
\noindent\textit{Editorial note.} The statement and proof are copied from the pinned chapter file; the cited inputs are kept as external pointers with their canonical Stacks tags.\par\medskip
\noindent\textit{Allowed standard prerequisites (editorial).}\begin{itemize}
\item Description of objects of the 2-fibre product. Same chapter, the description of the objects of a 2-fibre product of categories cofibred in groupoids as triples consisting of two objects and a morphism between their images. Supplies the translation of a diagram in the 2-fibre product into triples (u, v, a) and of morphisms into compatible triples.
\item (RS) property of the three categories. Same chapter, the statement of (RS) as the existence of fibre products over a fibre product of Artinian quotients. Supplies the fibre products u\_1 x\_u u\_2 and v\_1 x\_v v\_2 and, through the fibre square lemma, the identification of their images as fibre squares in the middle category.
\item Fibre squares in the middle category for images of fibre products. Same chapter, the lemma that the image under (RS) of a fibre product is again a fibre square. Supplies the two displayed fibre squares in the middle category, which allow the comparison morphism to be moved to the fibre product.
\end{itemize}\medskip
\ifdefined\hypertarget\hypertarget{sourcepointer-1}{}\fi
\noindent\textit{Source reference.} Original reference 06J3: sources/stacks-formal-defos/pinned-original-source.tex, source lines 5036--5036. Prerequisite pointer invoked at proof line(s) 5412; it names a standard fact of the same chapter that the counted argument uses, and every step of the counted proof built on it is retained verbatim in the delivered body.\par\medskip
\begin{lemma}
\label{lemma-RS-fiber-product-morphisms}
Let
$$
\xymatrix{
\mathcal{H} \times_\mathcal{F} \mathcal{G} \ar[r] \ar[d] &
\mathcal{G} \ar[d]^g \\
\mathcal{H} \ar[r]^f & \mathcal{F}
}
$$
be $2$-fibre product of categories cofibered in groupoids over
$\mathcal{C}_\Lambda$. If $\mathcal{F}, \mathcal{G}, \mathcal{H}$
all satisfy (RS), then $\mathcal{H} \times_\mathcal{F} \mathcal{G}$
satisfies (RS).
\end{lemma}

\begin{proof}
If $A$ is an object of $\mathcal{C}_\Lambda$, then an object of the fiber
category of $\mathcal{H} \times_\mathcal{F} \mathcal{G}$ over $A$
is a triple $(u, v, a)$ where $u \in \mathcal{H}(A)$, $v \in \mathcal{G}(A)$,
and $a : f(u) \to g(v)$ is a morphism in $\mathcal{F}(A)$. Consider
a diagram in $\mathcal{H} \times_\mathcal{F} \mathcal{G}$
$$
\vcenter{
\xymatrix{
           & (u_2, v_2, a_2) \ar[d] \\
(u_1, v_1, a_1) \ar[r] & (u, v, a)
}
}
\quad\text{lying over}\quad
\vcenter{
\xymatrix{
           & A_2 \ar[d] \\
A_1 \ar[r] & A
}
}
$$
in $\mathcal{C}_\Lambda$ with $A_2 \to A$ surjective. Since
$\mathcal{H}$ and $\mathcal{G}$ satisfy (RS), there are fiber
products $u_1 \times_u u_2$ and $v_1 \times_v v_2$ lying over
$A_1 \times_A A_2$. Since $\mathcal{F}$ satisfies (RS),
Lemma \ref{lemma-RS-fiber-square} shows
$$
\vcenter{
\xymatrix{
f(u_1 \times_u u_2) \ar[r] \ar[d] & f(u_2) \ar[d] \\
f(u_1) \ar[r] & f(u)
}
}
\quad\text{and}\quad
\vcenter{
\xymatrix{
g(v_1 \times_v v_2) \ar[r] \ar[d] & g(v_2) \ar[d] \\
g(v_1) \ar[r] & g(v)
}
}
$$
are both fiber squares in $\mathcal{F}$. Thus we can view
$a_1 \times_a a_2$ as a morphism from $f(u_1 \times_u u_2)$ to
$g(v_1 \times_v v_2)$ over $A_1 \times_A A_2$.
It follows that
$$
\xymatrix{
 (u_1 \times_u u_2, v_1 \times_v v_2, a_{1} \times_a a_2) \ar[d] \ar[r] &
(u_2, v_2, a_2) \ar[d] \\
(u_1, v_1, a_1) \ar[r] & (u, v, a)
}
$$
is a fiber square in $\mathcal{H} \times_\mathcal{F} \mathcal{G}$
as desired.
\end{proof}

\end{document}
